\begin{explanationbox}
\textbf{\textcolor{CExplanation}{先看图：弦中点与切线交点有什么关系}}
\begin{center}\begin{tikzpicture}[x=1cm,y=1cm,>=Stealth,font=\small]
\path[use as bounding box] (-2,-1.8) rectangle (10,6.4);
\node[anchor=west,font=\small\bfseries] at (-1.8,6) {斜弦的中点仍与 $P$ 同高，而 $P$ 不在弦上};
\draw[->,gray] (-1.6,0)--(6.6,0) node[right] {$x$};
\draw[->,gray] (0,-1.3)--(0,5.1) node[right] {$y$};
\draw[blue,thick,domain=-1.2:5.1,samples=81] plot ({\x*\x/4},\x);
\draw[orange!85!black,thick] (-1,2)--({3+2*sqrt(2)},{2+2*sqrt(2)});
\draw[orange!85!black,thick] (-1,2)--({3-2*sqrt(2)},{2-2*sqrt(2)});
\draw[red!75!black,very thick] ({3-2*sqrt(2)},{2-2*sqrt(2)})--({3+2*sqrt(2)},{2+2*sqrt(2)});
\draw[teal!80!black,dashed,thick] (-1.5,2)--(3.7,2);
\fill (-1,2) circle(2pt) node[above left] {$P(-1,2)$};
\fill (3,2) circle(2pt) node[below right] {$M(3,2)$};
\fill ({3+2*sqrt(2)},{2+2*sqrt(2)}) circle(2pt) node[above left] {$A$};
\fill ({3-2*sqrt(2)},{2-2*sqrt(2)}) circle(2pt) node[below left,xshift=-4pt] {$B$};
\node[anchor=north west,align=left,text width=3cm] at (7,4.8) {$y^2=4x$\\[6pt]红弦：$y=x-1$\\[10pt]$y_A=2+2\sqrt2$\\$y_B=2-2\sqrt2$\\[10pt]$y_M=\dfrac{y_A+y_B}{2}$\\[4pt]$\phantom{y_M}=2=y_0$};
\end{tikzpicture}\end{center}
图中 $P$ 是两条橙色切线的交点，$M$ 是红弦 $AB$ 的中点。虚线 $PM$ 水平，直接区分了“中点与 $P$ 同高”和“$P$ 在弦上”。

\textbf{\textcolor{CExplanation}{一、从已有知识出发：先认切线，再认切点弦}}
\begin{enumerate}[leftmargin=*,itemsep=0.4em]
\item $y^2=2px$ 上任一点 $Q(u,v)$ 满足 $v^2=2pu$；该点处切线为
\[\boxed{\ vy=p(x+u)\ }.\]
\item 若这条切线经过已知点 $P(x_0,y_0)$，那么 $vy_0=p(x_0+u)$。
\item 这个等式的特殊之处：\textbf{把 $Q(u,v)$ 看作变化的切点时，它的坐标必须满足同一个一次方程} $y_0y=p(x+x_0)$。两个切点自然落在同一条直线上。
\end{enumerate}


\medskip
\textbf{\textcolor{CExplanation}{二、从“一个点和一个斜率”求直线：相减与相加各做什么}}
这是一条\textbf{不需要先背切点弦公式}的理解路线：\textbf{直线 $AB$ 已知两个切点，就尝试求它的斜率；若用中点 $M$ 写点斜式，则只需找到中点坐标满足的等式。}

设 $A(x_1,y_1),B(x_2,y_2)$，其中点为 $M(x_m,y_m)$。两切点处的切线都经过 $P$，于是
\[
 \begin{cases}
 y_0y_1=p(x_0+x_1),\\
 y_0y_2=p(x_0+x_2).
 \end{cases}
\]
\textbf{相减：} 消掉公共的 $x_0$，只留下两切点坐标的差：
\[
 y_0(y_1-y_2)=p(x_1-x_2).
\]
当 $y_0\ne0$ 时，\(k_{AB}=\dfrac{y_1-y_2}{x_1-x_2}=\dfrac p{y_0}\)，\textbf{这是弦的方向}。

\textbf{相加：} 通过 $x_1+x_2=2x_m$、$y_1+y_2=2y_m$ 把两个端点联系到中点：
\[
 y_0y_m=p(x_0+x_m).
\]
\textbf{注意这只是中点坐标的一个关系，不是已经求出 $M$ 的具体坐标。}

\textbf{点斜式：} 以 $M$ 为已知点写直线：
\[
 y-y_m=\frac p{y_0}(x-x_m)
 \ \Longrightarrow\ y_0y=p(x+x_0).
\]
中点坐标虽然还没有单独算出，但代入时 $x_m,y_m$ 恰好可以消去。

\textbf{特殊情况：} 若 $y_0=0$，上面不能除以 $y_0$；两切线方程都给出 $x_1=x_2=-x_0$，所以弦为竖直线 $x=-x_0$，同样满足最终公式。\textbf{“相减求方向，相加找中点关系”要先判断斜率能否表示。}

\textbf{\textcolor{CExplanation}{三、为什么平方要变成“两坐标相乘”}}
抛物线方程中的 $y^2$ 只描述点在曲线上；而切线方程中的 $vy$ 同时包含\textbf{切点的纵坐标 $v$}与\textbf{直线上动点的纵坐标 $y$}。切点弦又把已知交点 $P$ 的纵坐标 $y_0$ 固定，因此出现 $y_0y$。

\medskip
\textbf{\textcolor{CExplanation}{四、为什么 $2px$ 变成 $p(x+x_0)$}}
当动点与切点重合，即 $(x,y)=(u,v)$ 时，切线公式自动变为 $v^2=p(u+u)=2pu$，恰好回到抛物线方程。这说明\textbf{系数必须是 $p$，横坐标必须相加}。

\medskip
\textbf{\textcolor{CExplanation}{五、判断存在条件：先算，再画}}
用切点纵坐标 $t$ 表示切点 $Q(t^2/(2p),t)$。因为其切线经过 $P$，所以
\[
 t y_0=p\left(x_0+\frac{t^2}{2p}\right)
 \iff t^2-2y_0t+2px_0=0.
\]
这是关于 $t$ 的一元二次方程：
\[
\Delta=(-2y_0)^2-4\cdot 1\cdot 2px_0
=4(y_0^2-2px_0).
\]
因此判断 $y_0^2-2px_0$ 的正负，恰好是在判断\textbf{能不能得到两个不同的实切点}。不能只凭口头所说的“外部点”来判断。

\medskip
\textbf{\textcolor{CExplanation}{六、从切点弦还能看出一个几何关系}}
设切点纵坐标是 $y_1,y_2$，上面的二次方程给出
\[
 y_1+y_2=2y_0,\qquad y_1y_2=2px_0.
\]
与图中两端点纵坐标的平均数对应，弦中点 $M$ 的纵坐标为 $\dfrac{y_1+y_2}{2}=y_0$，即\textbf{切点弦的中点与点 $P$ 在同一条水平线上}。再由 $x_i=y_i^2/(2p)$，可算得
\[
 x_M=\frac{x_1+x_2}{2}=\frac{y_0^2}{p}-x_0.
\]
注意：\textbf{中点与 $P$ 同高，不代表 $P$ 在切点弦上}。

\medskip
\textbf{\textcolor{CExplanation}{七、考场使用顺序}}
\begin{enumerate}[leftmargin=*,itemsep=0.3em]
\item 认开口：$y^2=2px$ 还是 $x^2=2py$。
\item 算正负：$y_0^2-2px_0$ 或 $x_0^2-2py_0$。
\item 若大于零，直接写切点弦方程；需要切点时再与抛物线联立。
\end{enumerate}
\end{explanationbox}
